%% kapa: exam=1 round=37 year=2026 period=1 form=A subject=civil_law subject_ko=민법 seq=28 number=28
%% source: Q-Net 감정평가사 기출문제 (cst00309) artlSeq=5258924 file=raw/1008/5258924/2026년 제37회 감정평가사 1차 1교시.pdf
%% conversion: kordoc 4.24.51 -> tools/gen1.py (auto); review=proofread-sonnet; flags=none
\begin{question}{28}
상린관계에 관한 설명으로 옳지 않은 것은? (다툼이 있으면 판례에 따름)
\begin{choices}
  \item 이웃 거주자는 생활방해가 수인한도를 초과하는 경우, 그 침해자에게 방해사유의 제거 내지 예방조치를 할 것을 청구할 수 있다.
  \item 토지소유자는 타인의 토지를 통과하지 아니하면 필요한 수도관 등을 시설할 수 없거나 과다한 비용을 요하는 경우, 그 타인의 동의없이 타인의 토지를 통과하여 이를 시설할 수 있다.
  \item 인지소유자는 자기의 비용으로 담의 재료를 통상보다 양호한 것으로 할 수 있다.
  \item 건물을 축조함에는 특별한 관습이 없으면 경계로부터 건물의 가장 돌출된 부분까지 반미터 이상의 거리를 두어야 한다.
  \item 인접지의 수목의 가지가 경계를 넘은 경우, 경계를 침범당한 사람은 수목의 소유자에 대한 가지제거 청구 없이 임의로 그 가지를 제거할 수 있다.
\end{choices}
\end{question}
\answer{5}
